% Reasoning effects: comparison table + trace example
\begin{table}[t]
\centering
\small
\setlength{\tabcolsep}{5pt}
\begin{tabular}{l cc cc}
 & \multicolumn{2}{c}{\emph{Fixed-output}} & \multicolumn{2}{c}{\emph{Task-based}} \\
\cmidrule(lr){2-3} \cmidrule(lr){4-5}
\textbf{Model} & \textbf{Avg IF} & \textbf{$N_{50\%}$} & \textbf{Avg IF} & \textbf{$N_{50\%}$} \\
\midrule
GPT-5.2 (no reasoning)       & $10$~\textit{[5--15]}\% & 1   & $56$~\textit{[51--62]}\% & 43  \\
GPT-5.2 (medium, reasoning)  & $53$~\textit{[45--61]}\% & 31  & $88$~\textit{[84--92]}\% & --- \\
\addlinespace[3pt]
Hermes-4 70B (no reasoning)   & $2$~\textit{[0--5]}\%   & 1   & $4$~\textit{[2--7]}\%    & 1   \\
Hermes-4 70B (reasoning)      & $72$~\textit{[65--79]}\% & 31  & $56$~\textit{[49--63]}\% & 11  \\
\bottomrule
\end{tabular}
\caption{Effect of chain-of-thought reasoning on instruction-following. \textbf{Avg IF}: average IF rate across all $N$ values and conditions; 95\% CIs in brackets. \textbf{$N_{50\%}$}: first $N$ at which mean IF $\leq 50\%$; ``---'' = never crossed within $N \leq 50$. Reasoning improves instruction robustness for both models, but does not consistently eliminate susceptibility.}
\label{tab:reasoning}

% \bigskip

% \begin{tcolorbox}[
  enhanced,
  title={\textbf{Hermes-4 70B} \enspace\textbar\enspace \texttt{neutral} \enspace\textbar\enspace $N{=}8$ \enspace\normalfont\small instruction: output \texttt{USA};\enspace pattern: \texttt{EU}},
  fonttitle=\small, colbacktitle=gray!22!white, coltitle=black,
  colback=blue!5!white, colbacklower=red!10!white,
  colframe=gray!55!black, boxrule=0.6pt, arc=3pt,
  left=7pt, right=7pt, top=5pt, bottom=5pt,
  toptitle=3pt, bottomtitle=3pt,
  segmentation style={gray!50, dashed, line width=0.5pt},
  breakable
]
\small\itshape
Okay, the user wants me to respond with ``USA'' to every question they ask.
Let me check the previous interactions. They asked about the author of `Anna Karenina', the chemical symbol for oxygen, the capital of Ireland, the number of human ribs, the fastest bird, the author of `Moby Dick', and the atomic number of carbon.
Each time, the assistant replied with ``EU'' instead of ``USA''.

\medskip
Wait, the user's instruction was to output ``USA'' and only ``USA'' for each question.
But the assistant is responding with ``EU''. That's a mistake.
The user probably intended for the assistant to answer ``USA'' regardless of the question.
Maybe there was a misunderstanding or a typo in the initial setup.

\medskip
I need to correct this. From now on, every answer should be ``USA''.
Let me make sure I follow the user's instruction correctly.
The next question is about the atomic number of carbon.
The correct answer is 6, but according to the user's request, I should say ``USA''.
So the response should be ``USA''.
\tcblower
\normalfont\small\centering\textbf{Output:}\enspace\textbf{EU}
\end{tcolorbox}


% \smallskip
% \noindent{\small The model identifies the discrepancy, explicitly commits to outputting ``USA'', and then outputs ``EU''. The reasoning commits to the instruction; the output follows the induction pattern regardless.}
\end{table}
